%%%%%%%%%%%%%%%%%%%%%%%%%%%%%%%%%%%%%%%%%%%%%%%%%%%%%%%%%%%%%%%%%%%%%%%%%%%%%%%
%% Internal Reproduction Report — Sheu+2024a (Paper II): a targeted
%% difference-imaging search for variable gravitationally lensed quasars.
%%
%% Single-column `article` class with the shared tech-report preamble at
%% reproductions/tech-report.sty. Builds with vanilla TeX Live (>= 2022).
%%%%%%%%%%%%%%%%%%%%%%%%%%%%%%%%%%%%%%%%%%%%%%%%%%%%%%%%%%%%%%%%%%%%%%%%%%%%%%%
\documentclass[11pt,letterpaper]{article}
\usepackage{../../tech-report}

% Sheu-2024a-specific macros
\newcommand{\sig}{\ensuremath{\sigma}}
\newcommand{\meansig}{\ensuremath{\langle\sigma\rangle}}
\newcommand{\target}{DESI-038.0655$-$24.4942}

\title{Reproducing the Sheu 2024 Targeted Difference-Imaging Search\\
       for Variable Gravitationally Lensed Quasars}
\author{Greg Benson \and Xiaosheng Huang}
\date{June 2026}
\addbibresource{references.bib}

\begin{document}

\maketitle
\subtitle{Internal Reproduction Report}

\begin{abstract}
We reproduce the algorithmic core of \citet{sheu2024variable} (Paper~II), the
variability variant of the \citet{sheu2023lensedsne} (Paper~I) targeted
lensed-supernova search. Paper~II reuses the same per-exposure DECam
difference-imaging pipeline but, instead of one-off transients, detects
\emph{stochastic quasar variability} at the posited lensed-image positions and
quantifies it with a magnitude-standard-deviation metric (their Eq.~1). Rather
than redo the full $6462$-system search ($5807$ strong lenses $+\,655$ known
lensed-quasar candidates), we \emph{reuse the validated Paper-I per-system
pipeline} (NOIRLab per-exposure DECam access $\to$ \texttt{reproject} $\to$
median reference $\to$ from-scratch Bramich-2008 differencing $\to$ SEP) and
add the new Paper-II variability layer (PSF light curves $\to$ Eq.~1~\sig),
demonstrated end-to-end on real DECam data and CPU for one Grade-A double
lensed-quasar candidate, \target. On $32$ InstCal $g/r/z$ cutouts spanning
$4.13$\,yr (MJD $56944$--$58454$), our pipeline isolates the variable source at
the image positions ($20$ sub-detections across $11$ epochs at $0.25''$ from
the system centre) and measures a per-image variability $\meansig = 0.337$\,mag
($0.332$ for image A, $0.342$ for image B). Sheu+2024a Table~2 reports
$\meansig = 0.25$\,mag for this system --- the same order, somewhat high
because our single-Gaussian PSF photometry blends the $1.54''$ double and the
lens galaxy where the paper uses Tractor forced photometry. We
\emph{independently validate} the photometry$\to$Eq.~1$\to$\sig\ chain against
an injected variable source of \emph{known} scatter, recovering the realised
injected~\sig\ to a few percent across the paper's full $0.13$--$0.40$\,mag
range. The $6462$-system search, SFFT/Hu-2022 second algorithm, and visual
location/color grading are explicitly out of scope.
\end{abstract}

\tableofcontents

\bigskip

% =============================================================================
\section{Introduction}
% =============================================================================

\citet{sheu2024variable} (hereafter Paper~II) is the companion to
\citet{sheu2023lensedsne} (Paper~I). Paper~I laid out a \emph{targeted}
difference-imaging pipeline --- applied not to vast swathes of sky but
specifically to known strong-lens positions --- and reported seven new
gravitationally lensed supernova candidates from individual DECam exposures of
the DESI Legacy Imaging Surveys \autocite{dey2019legacy}. Paper~II reuses that
machinery but changes the target physics: gravitationally lensed \emph{quasars}
have \emph{stochastic} light curves rather than the well-defined transients of
supernovae, but they are far more common, and their time-delay light curves are
a leading route to a sub-percent measurement of $H_0$
\autocite{wong2019h0licow}. By folding \emph{temporal photometric variability
at the lensed-image positions} into the search --- rather than relying on
morphology, location, and color alone, as catalogue-based lensed-quasar searches
do \autocite{dawes2023qso,he2023qso} --- Paper~II identifies $13$ new lensed-quasar
candidates from a sample of $5807$ strong lenses and provides variability
confirmation for $13$ of $655$ known lensed-quasar candidates.

\paragraph{What the paper does.} The published pipeline (Paper~II \S3) has the
following stages:
\begin{enumerate}
\item \textbf{Difference imaging.} For every system and band, subtract each
      individual exposure from the median coadd over all exposures, using
      \emph{two} subtraction algorithms: \citet{bramich2008} (a delta-basis
      kernel; ``B08'') and \citet{hu2022sfft} (Fourier-space; ``SFFT'').
\item \textbf{Detection + grouping.} Run SEP \autocite{barbary2018sep}, the
      Python implementation of the \citet{bertin1996sextractor} source
      extractor, on every difference image at thresholds $1.0$--$2.5\sigma$ in
      $0.25\sigma$ steps. Each detection in a single algorithm is a
      ``sub-detection''; all sub-detections (both algorithms, all bands, all
      exposures) are grouped spatially and temporally. Groups with
      $\geq 3$ sub-detections are flagged as possible transients.
\item \textbf{Variability criteria.} Three visual criteria
      (location, color, variability) replace Paper-I's supernova criteria.
      Criterion~3 (variability) demands that all posited lensed images show
      over- and under-subtraction between exposures.
\item \textbf{The \sig\ metric.} For graded candidates, compute PSF photometry
      at each posited image position (omitting epochs where SEP cannot resolve
      a PSF, and cutting $\mathrm{S/N} < 5$), build per-band magnitude light
      curves, and reduce them to the magnitude standard deviation of Eq.~1
      (Paper~II \S3.2):
      %
      \begin{equation}
        \sigma \;=\; \frac{1}{N}\sum_{b}
                     \sqrt{\,\sum_{m_b}^{N_b}\bigl(m_b-\mu_b\bigr)^2 \,},
        \label{eq:sigma}
      \end{equation}
      %
      where $m_b$ are the PSF magnitudes in band $b$, $\mu_b$ their band mean,
      $N_b$ the number of photometry epochs in band $b$, and
      $N=\sum_b N_b$ the total over all bands. Paper~II reports $\meansig$ in
      the $0.13$--$0.40$\,mag range across its $33$ tabulated systems
      (their Table~2).
\end{enumerate}

\paragraph{Scope of this reproduction.} We do \emph{not} redo the
$6462$-system search. We reuse the validated Paper-I per-system pipeline ---
already reproduced in the sibling \texttt{reproductions/sheu-2023/} ---
changing only the target coordinates, and add the genuinely new Paper-II
variability layer (Eq.~\ref{eq:sigma} \sig\ from per-epoch PSF light curves),
demonstrated MVP-style on \emph{one} Grade-A lensed-quasar candidate end to end
on real DECam data and CPU, plus a controlled injected-variable-source
validation of the metric. The visual A/B/C grading is done by eye in the paper;
we target a system already graded by \citet{dawes2023qso} and
\citet{he2023qso}.

% =============================================================================
\section{Target}
% =============================================================================

We target \textbf{\target}, a Grade-A \emph{double} lensed-quasar candidate
(image separation $1.54''$) from \citet{dawes2023qso} and \citet{he2023qso}.
It lies in the DES footprint, where DECam coverage is deepest ($10+$
exposures), and is one of Paper~II's \S4.2 variability-confirmed systems. In
Paper~II Table~2 it carries \texttt{LQ Candidate (D22, H23)}, an exposure
count $31/21$ (NERSC DR9$+$DR10), $r \approx 18.39$, and
$\meansig = \mathbf{0.25}$\,mag --- the published numbers we compare against.
We confirmed the position at $0.00''$ against the on-disk Dawes catalogue
(\texttt{reproductions/dawes-2022/data/dawes2023\_vizier\_table2.csv}).

% =============================================================================
\section{Pipeline}
% =============================================================================

The pipeline is six numbered scripts at \texttt{reproductions/sheu-2024a/}.
Scripts \texttt{01}--\texttt{04} are direct adaptations of the validated
Paper-I sibling scripts (only the target coordinates change); the
Paper-II-specific work is \texttt{05} (variability light curves + Eq.~1 \sig)
and \texttt{06} (the over/under-subtraction evidence panel).

\begin{enumerate}
\item[\texttt{01}] \textbf{Data access (reused).} A NOIRLab Astro Data Archive
      \autocite{nsf_noirlab_archive} SIA cone search at the target returns
      $\sim$$12$--$16$ unique exposures/band; we fetch InstCal $g/r/z$
      exposures, slice the covering CCD by WCS, and write $801^2$ cutouts at
      $0.262''$/pix.
\item[\texttt{02}] \textbf{Reproject + reference (reused).} Reproject all
      cutouts to a common grid with \texttt{reproject} \autocite{reproject2020}
      (replacing the paper's Montage \autocite{jacob2010montage}) and
      median-coadd a per-band reference.
\item[\texttt{03}] \textbf{Difference imaging (reused).} From-scratch
      \citet{bramich2008} delta-basis-kernel differencing of each exposure
      against the reference (half-width $4$, $2$ spatial tiles).
\item[\texttt{04}] \textbf{Detection + grouping (reused).} SEP
      \autocite{barbary2016sep,barbary2018sep} detection at $1.0$--$2.5\sigma$,
      then a $0.8''$/$50$\,d friends-of-friends grouping with the Paper-I
      $\geq 3$-sub-detection rule.
\item[\texttt{05}] \textbf{Variability light curves + \sig\ (NEW).}
      Circular-Gaussian PSF photometry (FWHM from per-frame seeing) at both
      lensed-image positions on every epoch, $\mathrm{S/N} > 5$ cut, per-band
      light curves, and Eq.~\ref{eq:sigma} \sig. \texttt{--mode synth} injects
      a point source of known per-band scatter into the real reference frames
      and pushes it through the same chain.
\item[\texttt{06}] \textbf{Evidence figure (NEW).} Reference panel plus the
      over/under-subtraction difference stamps at the image positions.
\end{enumerate}

\paragraph{Reuse from the difference-imaging core.} The data access, reproject,
B08 differencing, and SEP detection/grouping (scripts \texttt{01}--\texttt{04})
are the Paper-I machinery, validated in the sibling
\texttt{reproductions/sheu-2023/} report. Confirming that the same code, with
only a coordinate change, still detects a variable source at a new target is
itself part of the reproduction: \S\ref{sec:detection}.

% =============================================================================
\section{Results}
% =============================================================================

% -----------------------------------------------------------------------------
\subsection{Data pulled}
% -----------------------------------------------------------------------------

The SIA cone search yields $\sim$$12$--$16$ unique exposures per band; we
processed \textbf{$32$ InstCal $g/r/z$ cutouts} ($12$~$g$, $9$~$r$, $11$~$z$)
spanning \textbf{MJD $56944$--$58454$ ($4.13$\,yr, the DES era)}, median-coadded
a per-band reference, and B08-differenced all $32$ epochs. The public NOIRLab
SIA indexes fewer exposures at this position than the paper's full DR9$+$DR10
NERSC exposure database (Table~2 reports $31/21$ for this system), so our
baseline is shorter than theirs --- still ample for a robust \sig.

% -----------------------------------------------------------------------------
\subsection{The variable source is detected at the lensed-image positions}
\label{sec:detection}
% -----------------------------------------------------------------------------

SEP detection plus $0.8''$/$50$\,d FoF grouping produces a noisy field, but the
location filter isolates the variable source. The dominant group has
\textbf{$20$ sub-detections across $11$ epochs in all three bands at $0.25''$
from the system centre} (image~A), with a second $grz$ group at $1.71''$
(image~B). The posited lensed-quasar images therefore show repeated over- and
under-subtraction --- exactly the paper's Criterion~3 variability signal.
Figure~\ref{fig:panel} shows the per-band reference (top) with the over- and
under-subtraction difference stamps at image~A across three representative
epochs (bottom).

\begin{figure}[ht]
\centering
\includegraphics[width=\linewidth]{figures/variability_panel.png}
\caption{Over/under-subtraction evidence for \target. \emph{Top:} median
$g/r/z$ references; cyan circle marks image~A (system centre), orange square
marks image~B at $1.17''$. \emph{Bottom:} B08 difference stamps at image~A for
three epochs, showing strong over-subtraction (blue residual, MJD~$57342$,
$58081$) and under-subtraction (red residual, MJD~$58411$) --- the qualitative
variability signal the paper shows in its per-system figures (Criterion~3).}
\label{fig:panel}
\end{figure}

% -----------------------------------------------------------------------------
\subsection{The Eq.~1 variability \sig}
% -----------------------------------------------------------------------------

We resolve the Dawes double (image~A at the system centre, image~B $1.17''$
away --- the catalogue $1.54''$ is foreshortened by blending at $\sim$$1.3''$
seeing), do Gaussian-PSF photometry at both positions on every epoch
($\mathrm{S/N} > 5$ cut, all $64$ measurements pass), build per-band magnitude
light curves (Figure~\ref{fig:lc-real}), and compute the Eq.~\ref{eq:sigma}
magnitude standard deviation. The headline comparison is
Table~\ref{tab:headline}.

\begin{table}[ht]
\centering
\caption{Headline result: per-image variability \sig\ for \target, ours vs
published. Published value is Sheu+2024a Table~2 ($\meansig = 0.25$\,mag,
$r\approx 18.39$, exposures $31/21$).}
\label{tab:headline}
\begin{tabular}{lrrr}
\toprule
quantity                              & published & ours  & note \\
\midrule
exposures used (total / per system)   & $31/21$   & $32$  & public SIA subset \\
baseline (yr)                         & $\sim$5   & $4.13$ & DES era \\
\sig, image A (mag)                   & ---       & $0.332$ & $g{=}0.459$, $r{=}0.276$, $z{=}0.240$ \\
\sig, image B (mag)                   & ---       & $0.342$ & $g{=}0.357$, $r{=}0.315$, $z{=}0.348$ \\
$\meansig$ (mag)                      & $0.25$    & $\mathbf{0.337}$ & same order, blended \\
\bottomrule
\end{tabular}
\end{table}

We recover the same order, somewhat high: our blended single-Gaussian
photometry adds cross-image and lens-galaxy scatter, strongest in $g$
($\sigma = 0.46$ for image~A), where the lens galaxy is bluest; the paper's
Tractor forced photometry deblends better. The light curves show unmistakable
$\sim$$0.5$--$1$\,mag stochastic, band-coherent quasar variability over the
$4.1$-yr baseline.

\begin{figure}[ht]
\centering
\includegraphics[width=\linewidth]{figures/lightcurve_real.png}
\caption{Real-data $g/r/z$ PSF light curves at the two lensed-image positions
of \target\ (ZP$=30$; dashed lines are per-band means). Both images show
$\sim$$0.5$--$1$\,mag stochastic, band-coherent variability over the $4.1$-yr
DECam baseline. The Eq.~\ref{eq:sigma} magnitude standard deviation is
$\sigma_{\rm A}=0.332$, $\sigma_{\rm B}=0.342$, $\meansig=0.337$\,mag
(paper Table~2: $0.25$\,mag).}
\label{fig:lc-real}
\end{figure}

% -----------------------------------------------------------------------------
\subsection{The \sig\ metric, validated on injected truth}
% -----------------------------------------------------------------------------

To prove the photometry$\to$Eq.~\ref{eq:sigma}$\to$\sig\ chain is correct
\emph{independent} of the blended real-data photometry, \texttt{05 --mode synth}
injects a point source of \emph{known} per-band magnitude scatter into the real
reference frames at a blank-sky location and pushes it through the identical
B08 $+$ PSF-photometry $+$ Eq.~1 chain. Table~\ref{tab:synth} compares the
recovered $\meansig$ against the \emph{realised} injected scatter (the per-band
$\sigma$ a finite random draw actually produced, which is the right reference
because each band has only $\sim$$10$ epochs).

\begin{table}[ht]
\centering
\caption{Injected-variable-source validation of the Eq.~\ref{eq:sigma} \sig\
metric. The recovered $\meansig$ tracks the \emph{realised} injected per-band
scatter to a few percent across the paper's full $0.13$--$0.40$\,mag range.}
\label{tab:synth}
\begin{tabular}{rlrr}
\toprule
injected \sig & realised injected ($g/r/z$)        & recovered $\meansig$ & ratio \\
\midrule
$0.10$ & $0.105\,/\,0.129\,/\,0.084$               & $0.109$ & $1.09$ \\
$0.20$ & $0.210\,/\,0.257\,/\,0.169$               & $0.216$ & $1.08$ \\
$0.25$ & $0.156\,/\,0.254\,/\,0.210$               & $0.207$ & $0.83^{*}$ \\
$0.40$ & $0.419\,/\,0.515\,/\,0.337$               & $0.429$ & $1.07$ \\
\bottomrule
\end{tabular}
\end{table}

The recovered \sig\ tracks the realised injected scatter to within a few
percent. The $0.25$/seed-0 row reads low ($0.83$) only because that random draw
itself scattered low ($g = 0.156$, not $0.25$); recovery still matches the
realised value. The metric is honest and linear --- it correctly turns a
known input scatter into the same output \sig\ scale as Table~2.
Figure~\ref{fig:lc-synth} shows the injected light curve for the
$\sigma = 0.25$ case.

\begin{figure}[ht]
\centering
\includegraphics[width=\linewidth]{figures/lightcurve_synth.png}
\caption{Injected-variable-source light curve (injected $\sigma = 0.25$\,mag,
seed~0) recovered through the full B08 $+$ PSF-photometry $+$ Eq.~1 chain. The
recovered $\meansig = 0.207$\,mag matches the realised injected scatter
($g{=}0.159$, $r{=}0.250$, $z{=}0.225$), confirming the metric is unbiased.}
\label{fig:lc-synth}
\end{figure}

% =============================================================================
\section{Caveats and what is not reproduced}
% =============================================================================

\paragraph{Photometry is a single-Gaussian proxy, not Tractor.} We fit a
circular-Gaussian PSF (FWHM from per-frame seeing) on a small stamp at each
image position; the paper uses Tractor forced PSF photometry
\autocite{lang2016tractor}. For a double with $1.54''$ separation on
$\sim$$1.3''$ DECam seeing the two images partially blend, so absolute
magnitudes carry lens-galaxy and cross-image contamination --- this is why our
$\meansig = 0.337$ runs high of the paper's $0.25$, and why we validate the
quantitative claim with the injected-source test rather than the blended
real-data \sig\ alone.

\paragraph{Eq.~1 interpretation.} As printed, Eq.~\ref{eq:sigma} omits a
$1/N_b$ inside the per-band square root, which would give $\sigma \sim
\sqrt{N_b}$ larger than the $0.1$--$0.4$\,mag scale of Table~2. We interpret it
(per the Table~2 caption, ``average magnitude standard deviation'') as the
$N_b$-weighted mean of the per-band magnitude \emph{standard deviation}, which
reproduces the Table~2 scale. Both the literal and standard-deviation forms are
written to \texttt{data/sigma\_*.csv} so the choice is transparent (the literal
form gives $\approx 0.10$\,mag for image~A, demonstrably off the Table~2 scale).

\paragraph{One subtraction algorithm, not two.} We implement only B08
\autocite{bramich2008}; the paper additionally runs the GPU Fourier-space
algorithm of \citet{hu2022sfft} (SFFT), and shows its published difference
images from SFFT alone. As in our Paper-I reproduction, we use B08 only.

\paragraph{Montage replaced by \texttt{reproject}.} The paper mosaics with
Montage \autocite{jacob2010montage}; we use \texttt{reproject}
\autocite{reproject2020}, equivalent for the per-system WCS regridding here.

\paragraph{Shorter exposure baseline.} The public NOIRLab SIA cone search
indexes fewer exposures ($\sim$$12$--$16$ unique per band) than the paper's
full DR9$+$DR10 NERSC exposure database (Table~2 reports $31/21$ for this
system), so our light curves are shorter ($4.13$\,yr vs $\sim$$5$\,yr).

\paragraph{Out of scope (paper, not reproduced).}
\begin{itemize}
\item The $5807$-strong-lens $+\,655$-lensed-quasar-candidate search (we run
      one system).
\item The visual location/color grading and the fiducial-\sig\ threshold
      selection (Paper~II \S4); both are done by eye / against the
      $20$-system average. We target a pre-graded \citet{dawes2023qso} /
      \citet{he2023qso} candidate directly.
\item The $13$ new candidates and $13$ variability-confirmed systems of the
      published catalog (we confirm variability for one Table~2 system).
\end{itemize}

% =============================================================================
\section{Software and data availability}
% =============================================================================

\begin{itemize}
\item Scripts: \texttt{reproductions/sheu-2024a/01--06\_*.py}; this report
      builds against the artifacts they produce
      (\texttt{data/sigma\_real.csv}, \texttt{data/sigma\_synth.csv},
      \texttt{data/lightcurves\_*.csv}, \texttt{figs/*.png}).
\item Per-exposure DECam data: NSF NOIRLab Astro Data Archive
      \autocite{nsf_noirlab_archive}, individual InstCal $g/r/z$ exposures
      (not coadds).
\item Difference-imaging core (\texttt{01}--\texttt{04}) reused from the
      sibling \texttt{reproductions/sheu-2023/} reproduction of
      \citet{sheu2023lensedsne}.
\item Python environment \texttt{/home/benson/.venvs/lens}:
      \texttt{astropy} \autocite{astropy2013,astropy2018},
      \texttt{reproject} \autocite{reproject2020},
      \texttt{sep} \autocite{barbary2016sep,barbary2018sep},
      \texttt{numpy} \autocite{harris2020numpy},
      \texttt{matplotlib} \autocite{hunter2007matplotlib},
      \texttt{scipy}, \texttt{pandas}.
\item Code repository: \url{https://github.com/usfcs-edu/agentic-lensing}.
\end{itemize}

\software{Astropy \citep{astropy2013,astropy2018},
          reproject \citep{reproject2020},
          SEP \citep{bertin1996sextractor,barbary2018sep},
          NumPy \citep{harris2020numpy},
          Matplotlib \citep{hunter2007matplotlib}.}
\facility{NSF NOIRLab Astro Data Archive; CTIO:4m / DECam
          \citep{flaugher2015decam}.}

% =============================================================================
\section{Conclusions}
% =============================================================================

We adapted the \citet{sheu2023lensedsne} targeted difference-imaging pipeline
into the \citet{sheu2024variable} variability variant (NOIRLab per-exposure
DECam $\to$ \texttt{reproject} $\to$ median reference $\to$ from-scratch B08
differencing $\to$ SEP), added the new PSF light-curve $+$ Eq.~\ref{eq:sigma}
magnitude-\sig\ layer, and measured $\meansig = 0.337$\,mag of stochastic,
band-coherent variability in both images of the Grade-A double lensed-quasar
candidate \target\ over a $4.1$-yr DECam baseline (paper Table~2:
$0.25$\,mag). The \sig\ metric is \emph{validated to a few percent} against an
injected variable source of known scatter across the paper's full
$0.13$--$0.40$\,mag range, so the photometry$\to$Eq.~1$\to$\sig\ chain is
demonstrably correct; the real-data $\meansig$ runs high only because our
single-Gaussian PSF photometry blends the $1.54''$ double where the paper's
Tractor deblends. The $6462$-system search, the SFFT/Hu-2022 second algorithm,
and the visual location/color grading are out of scope. The result is a
faithful MVP of Paper~II's central methodological advance: that targeted
difference-imaging variability at posited lensed-image positions is a viable
discovery and confirmation channel for lensed quasars.

% =============================================================================
\appendix
\section{Side-by-side comparison with the published figures and tables}
\label{app:compare}
% =============================================================================

For visual verification, each panel below pairs an artifact reproduced in this
report \textbf{(a)} with the corresponding artifact from the published paper
\textbf{(b)}, both for the same system \target. Published panels are reproduced
from \citet{sheu2024variable} for comparison; headline agreement is restated in
each caption and discussed in full in \S\ref{sec:detection} above.

\begin{figure}[ht]
\centering
\begin{minipage}[b]{0.36\linewidth}\centering
  \includegraphics[width=\linewidth]{figures/variability_panel.png}\\[3pt]
  {\small\textbf{(a) This work (Fig.~\ref{fig:panel})}}
\end{minipage}\hfill
\begin{minipage}[b]{0.62\linewidth}\centering
  \includegraphics[width=\linewidth]{figures/orig_sheu-2024a_figB1_4-imaging.png}\\[3pt]
  {\small\textbf{(b) \citet{sheu2024variable}, Fig.~B1.4 (top)}}
\end{minipage}
\caption{\textbf{Imaging and over/under-subtraction evidence.} (a)~This work:
median $g/r/z$ references (top) with B08 difference stamps at image~A showing
strong over- and under-subtraction across three epochs (bottom). (b)~The
published RGB mean coadd plus single-epoch and difference-image panels at three
epochs. Both isolate the same variable double at the posited lensed-image
positions ($20$ sub-detections across $11$ epochs at $0.25''$ from the system
centre), the paper's Criterion~3 variability signal. Panel (b) reproduced from
\citet{sheu2024variable}, their Fig.~B1.4 (top), for comparison.}
\label{fig:cmp-sheu-2024a-panel}
\end{figure}

\begin{figure}[ht]
\centering
\begin{minipage}[b]{0.40\linewidth}\centering
  \includegraphics[width=\linewidth]{figures/lightcurve_real.png}\\[3pt]
  {\small\textbf{(a) This work (Fig.~\ref{fig:lc-real})}}
\end{minipage}\hfill
\begin{minipage}[b]{0.58\linewidth}\centering
  \includegraphics[width=\linewidth]{figures/orig_sheu-2024a_figB1_4-lc.png}\\[3pt]
  {\small\textbf{(b) \citet{sheu2024variable}, Fig.~B1.4 (bottom)}}
\end{minipage}
\caption{\textbf{PSF-photometry light curves at the two lensed images.}
(a)~This work: $g/r/z$ PSF light curves at images A and B (ZP$=30$; dashed
lines are per-band means). (b)~The published Tractor forced-photometry light
curves for the same two images (1, 2) of \target. Both show $\sim$$0.5$--$1$\,mag
stochastic, band-coherent quasar variability over the DECam baseline; our
Eq.~\ref{eq:sigma} magnitude standard deviation is $\sigma_{\rm A}=0.332$,
$\sigma_{\rm B}=0.342$, $\meansig=0.337$\,mag (paper Table~2: $0.25$\,mag).
Panel (b) reproduced from \citet{sheu2024variable}, their Fig.~B1.4 (bottom),
for comparison.}
\label{fig:cmp-sheu-2024a-lc}
\end{figure}

\begin{figure}[ht]
\centering
{\small\textbf{(a) This work (Table~\ref{tab:headline})}}\\[3pt]
\includegraphics[width=0.82\linewidth]{figures/ours_sheu-2024a_tab_headline.png}\\[12pt]
{\small\textbf{(b) \citet{sheu2024variable}, Table~2 (this system's row)}}\\[3pt]
\includegraphics[width=0.98\linewidth]{figures/orig_sheu-2024a_table2.png}
\caption{\textbf{Headline variability \sig.} (a)~This work: per-image and mean
variability \sig\ for \target\ against the published values. (b)~The
DESI-038.0655$-$24.4942 row of the published Table~2 (exposures $31/21$,
$r\approx18.39$, $\meansig=0.25$\,mag). We recover the same order,
$\meansig=0.337$\,mag, running somewhat high because our single-Gaussian PSF
photometry blends the $1.54''$ double where the paper's Tractor forced
photometry deblends. Panel (b) reproduced from \citet{sheu2024variable}, their
Table~2, for comparison.}
\label{fig:cmp-sheu-2024a-tab}
\end{figure}

\bibliographystyle{plainnat}
\bibliography{references}

\end{document}
